% Replace the active Figure 2.1 environment with the one below. Keep its
% existing \label so older cross-references continue to work. The paragraph
% and scope table can follow at the start of the study-specific discussion.
\begin{figure}[htbp]
  \centering
  \caption{From welding evidence to requirements for testing fusion}
  \label{fig:fusion-research-gap}
  \includegraphics[width=\linewidth]{figures/fusion_research_gap.pdf}
  \displaynote{DMAW reported 96.1\% accuracy for fusion and 90.2\% for
    its strongest single input \parencite{P003}. In weld penetration, both
    single-input models and fusion reached 100\% on the clean test set
    \parencite{Hua2026Weld}. The reported contrasts are within each study.}
\end{figure}

The comparison starts with a valid target and observations paired to the same
event (R1). For each input, the choice of representation and fitting budget
defines the single-input reference (R2, SQ1). Fusion is then compared with the
strongest reference on the same cases, with matched opportunities for
recalibration (R3, SQ2). In the initial full-text review, four of 37 papers
examined planned missing-input inference and one examined simulated dropout
(13.5\% combined, Table~\ref{tab:literature-matrix}). Reliability comparisons
therefore distinguish absence, corruption, and distribution shift. When testing
adaptive weighting, the static comparator needs equal training exposure to
the tested failures (R4, SQ3). Grouped splits, leakage checks, and uncertainty
estimates bound the claims. For classification, precision--recall analysis
and calibration address different aspects of performance (R5, SQ4).
Table~\ref{tab:fusion-study-scope} identifies the comparisons each study
could support.

\begin{table}[htbp]
  \centering
  \caption{Inputs and assessment limits of the five thesis studies}
  \label{tab:fusion-study-scope}
  \begingroup
  \singlespacing\small
  \setlength{\tabcolsep}{4pt}
  \renewcommand{\arraystretch}{1.28}
  \begin{tabularx}{\linewidth}{@{}P{.22\linewidth} P{.23\linewidth} Y@{}}
    \toprule
    \textbf{Study} & \textbf{Predictor inputs} & \textbf{Scope and assessment limits}\\
    \midrule
    PHM (Ch.~5)
      & Machine sensors
      & R1--R3, R5. Post-run photographs define the target proxy.
        The study does not test camera outages.\\
    REASSEMBLE (Ch.~6)
      & Video and machine sensors
      & R1--R5. Principal study of fusion and of adaptive versus static
        weighting under matched corruption exposure.\\
    MSM (Ch.~7)
      & Video and machine telemetry
      & R1--R5. The comparison covers cutting intervals observed in the
        laboratory during acquisition.\\
    ViFailback (Ch.~8)
      & Video features and robot state
      & R1--R3, R5. The planned corruption stress test (R4) was unexecuted.\\
    HANDOVER (App.~\ref{app:handover})
      & Video and force--torque measurements
      & R1--R5 with fixed averaging. The missing participant-to-trial mapping
        precluded a participant-disjoint fit of the learned stacker.\\
    \bottomrule
  \end{tabularx}
  \endgroup
  \displaynote{The R codes indicate which requirements a study investigated.
    Each conclusion applies to that study's target and assessment population.}
\end{table}
